\documentclass[12pt]{article}
\usepackage{amsmath, amssymb, geometry, setspace}
\geometry{margin=1in}
\setstretch{1.2}
\setlength{\parindent}{0pt}



\title{Solutions -- Quiz 4\\ \vspace{10pt}
\large ECON 320 -- Introduction to Mathematical Economics}
\author{Instructor: Muhebullah Karimzada}
\date{October 22, 2025}
\begin{document}
\maketitle

\section*{Part A – Multiple Choice Questions (1 mark each)}

\begin{enumerate}
\item \textbf{(a)} 
\[
\frac{d}{dx}\ln(x^2+1)=\frac{1}{x^2+1}\cdot(2x)=\frac{2x}{x^2+1}.
\]

\item \textbf{(b)} 
\[
\frac{d}{dx}e^{2x}=2e^{2x}.
\]

\item \textbf{(b)} 
\[
y=e^{3x-5}\Rightarrow \frac{dy}{dx}=3e^{3x-5}.
\]

\item \textbf{(b)}  
\(\ln(x)\) is defined for \(x>0\).

\item \textbf{(c)}  
\[
f(x)=\ln(e^{2x})=2x \Rightarrow f'(x)=2.
\]

\item \textbf{(b)}  
For \(y=ae^{bx}\),
\[
E_{yx}=\frac{dy}{dx}\frac{x}{y}=abe^{bx}\frac{x}{ae^{bx}}=b.
\]

\item \textbf{(a)}  
\(\ln(x)\) is increasing (\(f'(x)=1/x>0\)) and concave (\(f''(x)=-1/x^2<0\)).

\end{enumerate}

\newpage

\section*{Part B – Analytical Problems (13 marks total)}

\subsection*{Question 1 (4 marks)}

Utility: \(U(x)=\ln(x+2)\).

\textbf{(a) Derivatives.}
\[
U'(x)=\frac{1}{x+2}, \quad U''(x)=-\frac{1}{(x+2)^2}.
\]

\textbf{(b) Monotonicity and Concavity.}
\[
U'(x)>0 \ \forall x>-2 \Rightarrow \text{Utility increasing.}
\]
\[
U''(x)<0 \ \forall x>-2 \Rightarrow \text{Utility concave.}
\]

\textbf{(c) Economic interpretation.}
Concavity indicates \textbf{diminishing marginal utility of income}: each extra dollar adds less to utility.

\bigskip

\subsection*{Question 2 (4 marks)}

Production: \(Q(L)=20L^{0.5}e^{-0.05L}\).

\textbf{(a) Derivative.}
\[
Q'(L)=20e^{-0.05L}(0.5L^{-0.5}-0.05L^{0.5})
=20e^{-0.05L}L^{-0.5}(0.5-0.05L).
\]

\textbf{(b) FOC.}
\[
0.5-0.05L=0 \Rightarrow L^*=10.
\]

\textbf{(c) SOC.}
\[
Q''(L)=20e^{-0.05L}L^{-1.5}(-0.025L^2+0.0025L-0.25),
\]
which is negative at \(L=10\), so \(L^*\) yields a maximum.

\textbf{(d) Interpretation.}  
The exponential term models diminishing productivity—beyond \(L^*\), crowding or inefficiency lowers output.

\bigskip
\newpage
\subsection*{Question 3 (5 marks)}

Demand: \(Q(P)=100e^{-0.2P}\).  
Revenue: \(R(P)=PQ(P)=100P e^{-0.2P}\).

\textbf{(a) Derivative.}
\[
R'(P)=100e^{-0.2P}(1-0.2P).
\]

\textbf{(b) FOC.}
\[
1-0.2P=0\Rightarrow P^*=5.
\]

\textbf{(c) Maximum revenue.}
\[
R(5)=100(5)e^{-1}=500/e\approx183.94.
\]

\textbf{(d) SOC.}
\[
R''(P)=100e^{-0.2P}(-0.2(1-0.2P)-0.2)
=100e^{-0.2P}(-0.4+0.04P).
\]
At \(P=5\): \(R''(5)=100e^{-1}(-0.2)<0\Rightarrow\) maximum.

\textbf{(e) Interpretation.}  
At low prices, higher prices increase total revenue; beyond \(P=5\), the exponential decay in demand dominates, reducing total revenue.



\end{document}
